\documentclass{article}
\usepackage[utf8]{inputenc}
\usepackage{amsmath,amssymb}
\usepackage[most]{tcolorbox}
\usepackage{geometry}
\geometry{margin=2.5cm}

\newcommand{\step}[1]{\par\noindent\hangindent=3.5em\hangafter=1%
  \makebox[3.5em][l]{\textbf{#1.}}}
\newcommand{\steptag}[1]{\hfill{\small\textsf{[#1]}}}

\newtcolorbox{proofbox}[1]{
  colback=gray!5, colframe=gray!60,
  fonttitle=\bfseries, title={#1},
  breakable, enhanced,
  left=4pt, right=4pt, top=4pt, bottom=4pt,
  before skip=10pt, after skip=10pt
}

\begin{document}

\begin{proofbox}{After first fixing one global witness package W\_RF suppli...}
\step{1} After first fixing one global witness package W\_RF supplied by lem-routef-raw-factor-setting-formation, for every input (H,Phi,eta) to which that formation result applies, fix one def-routef-raw-factor-setting datum S over that same W\_RF supplied by the same result, for every Delta' supplied for (W\_RF,S) by lem-routef-delta-prime-closeness, and every Delta supplied from that same Delta' by lem-routef-delta-normalization-closeness; for every integer n >= 1 and all X, Y, Z in M\_n(S.B), writing the fields of (W\_RF,S) as the unqualified symbols below: Route F degree-three estimate: with C\_3 := 10+20*C\_Delta+12*C\_theta+2*C\_Delta' and rho\_3 := min\{rho\_theta, rho\_Delta', rho\_Delta, rho\_2\}, for 0 <= eta <= rho\_3, every amplification satisfies ||Phi\_n(Delta\_n X Delta\_n Y Delta\_n Z) - Delta\_n(XYZ)|| <= C\_3*eta*||X||*||Y||*||Z||. \steptag{validated} \\[6pt]
\step{1.1} Admissibility and the one-factor invariance estimate. Under the root hypotheses and 0 <= eta <= rho\_3, all cited estimates below are applicable, Phi and Delta are UCP and hence completely contractive, tilde-Phi\_n tilde-Delta\_n=tilde-Delta\_n, ||tilde-Delta||\_cb <= 2, and for every amplification n and T in M\_n(B), ||Phi\_n(Delta\_n T)-Delta\_n T|| <= d||T|| with d:=2(C\_theta+C\_Delta)eta. \steptag{validated} \\[6pt]
\step{1.1.1} Radius bookkeeping, map status, and exact range identity. By def-routef-raw-factor-setting, rho\_3=min\{rho\_theta,rho\_Delta\_prime,rho\_Delta,rho\_2\}, rho\_theta=1/8, and rho\_Delta\_prime<=rho\_T. Thus eta<=rho\_3 places us simultaneously in the domains of lem-routef-functional-calculus-closeness, lem-routef-raw-factor-norms, lem-routef-delta-prime-closeness, lem-routef-delta-normalization-closeness, lem-routef-degree-two-estimate, and in eta<1/4 for lem-kitaev-almost-idemp-audit. The root choice is licensed by lem-routef-raw-factor-setting-formation; Phi is UCP there, and the chosen Delta is UCP by lem-routef-delta-normalization-closeness, so both are completely contractive. Formation also gives tilde-Phi\textasciicircum{}2=tilde-Phi and tilde-Delta(B) subset A=Im(tilde-Phi); hence tilde-Phi\_n tilde-Delta\_n=tilde-Delta\_n at every amplification. \steptag{validated} \\[4pt]
\step{1.1.2} Uniform raw-map bound. Lem-routef-raw-factor-norms gives ||tilde-Delta||\_cb<=1+C\_T eta because eta<=rho\_3<=rho\_Delta\_prime<=rho\_T. From def-routef-raw-factor-setting, C\_T=C\_theta+3C\_V and rho\_T<=1/[4(1+C\_theta)] and rho\_T<=1/[4(1+C\_V)]. Hence C\_theta eta<=1/4 and 3C\_V eta<=3/4, so C\_T eta<=1 and ||tilde-Delta||\_cb<=2. \steptag{validated} \\[4pt]
\step{1.1.3} One-factor estimate. Using tilde-Phi\_n tilde-Delta\_n=tilde-Delta\_n, the exact identity Phi\_n Delta\_n-Delta\_n=(Phi\_n-tilde-Phi\_n)tilde-Delta\_n+Phi\_n(Delta\_n-tilde-Delta\_n)+(tilde-Delta\_n-Delta\_n) holds. Lem-routef-functional-calculus-closeness bounds the first map difference by C\_theta eta, lem-routef-delta-normalization-closeness bounds ||Delta-tilde-Delta||\_cb by C\_Delta eta, Phi is completely contractive, and ||tilde-Delta||\_cb<=2. Therefore for every T, ||Phi\_n(Delta\_n T)-Delta\_n T||<=[2C\_theta+C\_Delta+C\_Delta]eta||T||=2(C\_theta+C\_Delta)eta||T||=d||T||. \steptag{validated} \\[4pt]
\step{1.2} First three-factor replacement. Assuming the one-factor estimate ||Phi\_n(Delta\_n T)-Delta\_n T|| <= d||T|| and complete contractivity of Phi and Delta, put P0:=Phi\_n(Delta\_n X Delta\_n Y Delta\_n Z) and P1:=Phi\_n(Phi\_n(Delta\_n X) Phi\_n(Delta\_n Y) Phi\_n(Delta\_n Z)); a three-term product telescope gives ||P0-P1|| <= 3d||X||||Y||||Z||. \steptag{validated} \\[6pt]
\step{1.2.1} Write D\_X=Delta\_n X and F\_X=Phi\_n(D\_X), and similarly for Y,Z. The exact product identity D\_XD\_YD\_Z-F\_XF\_YF\_Z=(D\_X-F\_X)D\_YD\_Z+F\_X(D\_Y-F\_Y)D\_Z+F\_XF\_Y(D\_Z-F\_Z), followed by complete contractivity of Delta and Phi and the assumed one-factor estimate, bounds its norm by 3d||X||||Y||||Z||. Applying the contractive outer Phi\_n gives ||P0-P1||<=3d||X||||Y||||Z||. \steptag{validated} \\[4pt]
\step{1.3} Associativity comparison. Define P1:=Phi\_n(Phi\_n(Delta\_n X) Phi\_n(Delta\_n Y) Phi\_n(Delta\_n Z)) and P2:=Phi\_n(Phi\_n(Phi\_n(Delta\_n X) Phi\_n(Delta\_n Y)) Phi\_n(Delta\_n Z)). Since eta <= rho\_3 <= rho\_theta=1/8<1/4, the amplified Phi\_assoc1 estimate of lem-kitaev-almost-idemp-audit applied to Delta\_n X, Delta\_n Y, Delta\_n Z, together with complete contractivity of Delta, gives ||P1-P2|| <= 10eta||X||||Y||||Z||. \steptag{validated} \\[6pt]
\step{1.3.1} At amplification n, Phi\_assoc1 from lem-kitaev-almost-idemp-audit is precisely ||Phi\_n(Phi\_n(Phi\_n(U)Phi\_n(V))Phi\_n(W))-Phi\_n(Phi\_n(U)Phi\_n(V)Phi\_n(W))||<=10eta||U||||V||||W|| whenever eta<1/4. Substitute U=Delta\_n X, V=Delta\_n Y, W=Delta\_n Z. The two displayed terms are P2 and P1 respectively, while UCP contractivity gives ||U||||V||||W||<=||X||||Y||||Z||; symmetry of norm under sign reversal yields the claimed comparison. \steptag{validated} \\[4pt]
\step{1.4} Second three-factor replacement. Define P2:=Phi\_n(Phi\_n(Phi\_n(Delta\_n X) Phi\_n(Delta\_n Y)) Phi\_n(Delta\_n Z)) and P3:=Phi\_n(Phi\_n(Delta\_n X Delta\_n Y) Delta\_n Z). Assuming ||Phi\_n(Delta\_n T)-Delta\_n T|| <= d||T|| for all T, where d:=2(C\_theta+C\_Delta)eta, and complete contractivity of Phi and Delta, a two-factor telescope inside the inner Phi\_n costs at most 2d||X||||Y|| and replacing the last Phi\_n(Delta\_n Z) by Delta\_n Z costs at most d||Z||, so ||P2-P3|| <= 3d||X||||Y||||Z||. \steptag{validated} \\[6pt]
\step{1.4.1} Let A=Phi\_n(Phi\_n(Delta\_n X)Phi\_n(Delta\_n Y)), B=Phi\_n(Delta\_n Z), C=Phi\_n(Delta\_n X Delta\_n Y), and D=Delta\_n Z, so P2=Phi\_n(AB) and P3=Phi\_n(CD). Contractivity of the outer Phi\_n and AB-CD=(A-C)B+C(B-D) give ||P2-P3||<=||A-C||||B||+||C||||B-D||. The exact two-factor telescope Phi\_n(Delta\_n X)Phi\_n(Delta\_n Y)-Delta\_n X Delta\_n Y=[Phi\_n(Delta\_n X)-Delta\_n X]Phi\_n(Delta\_n Y)+Delta\_n X[Phi\_n(Delta\_n Y)-Delta\_n Y], followed by contractivity, gives ||A-C||<=2d||X||||Y||. Also ||B||<=||Z||, ||C||<=||X||||Y||, and ||B-D||<=d||Z||. Thus ||P2-P3||<=3d||X||||Y||||Z||. \steptag{validated} \\[4pt]
\step{1.5} Degree-two finish and coefficient closure. Define P3:=Phi\_n(Phi\_n(Delta\_n X Delta\_n Y) Delta\_n Z), P4:=Phi\_n(Delta\_n(XY) Delta\_n Z), and P5:=Delta\_n(XYZ). Because eta <= rho\_3 <= rho\_2, two applications of lem-routef-degree-two-estimate, first to X,Y and then to XY,Z, plus complete contractivity, give ||P3-P4|| <= C\_2 eta||X||||Y||||Z|| and ||P4-P5|| <= C\_2 eta||X||||Y||||Z||. Independently, for d:=2(C\_theta+C\_Delta)eta and C\_2=C\_Delta'+4C\_Delta, 6d+10eta+2C\_2 eta=(10+12C\_theta+20C\_Delta+2C\_Delta')eta=C\_3 eta. \steptag{validated} \\[6pt]
\step{1.5.1} First degree-two comparison. Contractivity of the outer Phi\_n gives ||P3-P4||<=||Phi\_n(Delta\_n X Delta\_n Y)-Delta\_n(XY)|| ||Delta\_n Z||. Lem-routef-degree-two-estimate applies because eta<=rho\_3<=rho\_2 and bounds the first factor by C\_2 eta||X||||Y||; UCP contractivity of Delta bounds the second by ||Z||. Hence ||P3-P4||<=C\_2 eta||X||||Y||||Z||. \steptag{validated} \\[4pt]
\step{1.5.2} Second degree-two comparison. Apply lem-routef-degree-two-estimate at amplification n to the pair XY,Z. Since M\_n(B) is an associative C*-algebra, ||XY||<=||X||||Y|| and (XY)Z=XYZ, so ||P4-P5||=||Phi\_n(Delta\_n(XY)Delta\_n Z)-Delta\_n((XY)Z)||<=C\_2 eta||XY||||Z||<=C\_2 eta||X||||Y||||Z||. \steptag{validated} \\[4pt]
\step{1.5.3} Coefficient arithmetic. By def-routef-raw-factor-setting, C\_2=C\_Delta\_prime+4C\_Delta and C\_3=10+20C\_Delta+12C\_theta+2C\_Delta\_prime. With d=2(C\_theta+C\_Delta)eta, 6d+10eta+2C\_2 eta=[12C\_theta+12C\_Delta+10+2C\_Delta\_prime+8C\_Delta]eta=C\_3 eta. \steptag{validated} \\[4pt]
\end{proofbox}

\end{document}
